%  LaTeX support: latex@mdpi.com 
%  For support, please attach all files needed for compiling as well as the log file, and specify your operating system, LaTeX version, and LaTeX editor.

%=================================================================
\documentclass[jmse,article,submit,pdftex,moreauthors]{Definitions/mdpi} 

%=================================================================
%----------
% journal

% MDPI internal commands
\firstpage{1} 
\makeatletter 
\setcounter{page}{\@firstpage} 
\makeatother
\pubvolume{1}
\issuenum{1}
\articlenumber{0}
\pubyear{2024}
\copyrightyear{2024}
%\externaleditor{Academic Editor: Firstname Lastname}
\datereceived{7 March 2024} 
\daterevised{}
\dateaccepted{} 
\datepublished{} 
\hreflink{https://doi.org/} % If needed use \linebreak
%\doinum{}
\usepackage{lipsum} 

%=================================================================
% Add packages and commands here. The following packages are loaded in our class file: fontenc, inputenc, calc, indentfirst, fancyhdr, graphicx, epstopdf, lastpage, ifthen, lineno, float, amsmath, setspace, enumitem, mathpazo, booktabs, titlesec, etoolbox, tabto, xcolor, soul, multirow, microtype, tikz, totcount, changepage, attrib, upgreek, cleveref, amsthm, hyphenat, natbib, hyperref, footmisc, url, geometry, newfloat, caption

\graphicspath{{figures/}} % path to folder with figures
\usepackage{subcaption}
\usepackage{listings}
\usepackage{xcolor}
\usepackage{gensymb} % for \textdegree
\usepackage{tabularx}

%=================================================================
% Full title of the paper (Capitalized)
\Title{Artificial Neural Networks for Mapping Coastal Lagoon of Chilika Lake, India, Using Earth Observation Data}

% MDPI internal command: Title for citation in the left column
\TitleCitation{Artificial Neural Networks for Mapping Coastal Lagoon of Chilika Lake, India, Using Earth Observation Data}

% Author Orchid ID: enter ID or remove command
\newcommand{\orcidauthorA}{0000-0002-5759-1089} % Polina Add \orcidA{} behind the author's name

% Authors, for the paper (add full first names)
\Author{Polina Lemenkova $^{1}$*\orcidA{}}

%\longauthorlist{yes}

% MDPI internal command: Authors, for metadata in PDF
\AuthorNames{Polina Lemenkova}

% MDPI internal command: Authors, for citation in the left column
\AuthorCitation{Lemenkova, P.}
% If this is a Chicago style journal: Lastname, Firstname, Firstname Lastname, and Firstname Lastname.

% Affiliations / Addresses (Add [1] after \address if there is only one affiliation.)
\address{%
$^{1}$ \quad Department of Geoinformatics, Faculty of Digital and Analytical Sciences, Universität Salzburg, Schillerstraße 30, A-5020 Salzburg, Austria.
}

% Contact information of the corresponding author
\corres{Correspondence: polina.lemenkova@plus.ac.at; Tel.: +43-677-6173-2772 (P.L.)}

% Current address and/or shared authorship
%\firstnote{Current address: Affiliation 3.} 

% Abstract (Do not insert blank lines, i.e. \\) 
\abstract{\textcolor{red}{\lipsum[1-2]}Functional algorithm of gradient boosting classifier }

% Keywords
\keyword{Functional algorithm; risk assessment; ensemble learning; hazards; climate change; GRASS GIS; Scikit-Learn; machine learning; Python; Landsat; image analysis; China; classification} 

% The fields PACS, MSC, and JEL may be left empty or commented out if not applicable
%\PACS{J0101}
%\MSC{}
%\JEL{}

\begin{document}

%----------- section ----------->
\section{Introduction}

\subsection{Background}
Coastal and brackish water lagoons that form continuities in the terrestrial and shelf ecosystems present the transitional aquatic ecosystems located in the zones between land and sea. They are often associated with dynamic environmental conditions and high biodiversity due to their connections to the marine and terrestrial communities \cite{PEREZRUZAFA2019253}. Formed as a result of the marine transgression, coastal lagoons vary with depth and geomorphic features, tidal circulation patterns, salinity and wind forcing \cite{KJERFVE198663}. The formation of the barriers between the land and ocean is driven by the controversial forces of erosion and sediment deposition within the coastal lagoons \cite{BOYD1992139}. The accumulation of river sediments depends on the river speed which transport sediments in large quantity, and intensified through complex hydrological processes such as winds, oceanic currents and coastal waves \cite{KUMAR2016155}, and groundwater discharge \cite{DANISH2020103816}. Moreover, different slope and substrate types and density-driven currents with diverse morphodynamic tidal regime affect littoral zones of the coastal lagoon and create stratified conditions of the sediment resuspension \cite{BORTOLIN2020105782,Larson2012}. 

Topographically isolated and hydrologically distinct from the surrounding landscapes, coastal lagoons include unique species of wild flora and fauna (e.g., bird species). Border situation of coastal lagoons determines their high levels of biodiversity through the intense physical-chemical gradients. The particular feature of coastal lagoons which ensures its high productivity as aquatic systems consists in shallow bathymetry \cite{PEREZRUZAFA2024171264}. Since sunlight reaches all levels the shallow  systems of the coastal lagoons until the lowest bottom layer and the water due to the active wave and currents, the water mass of the lagoon is well mixed. Such dynamics activates the recycling of nutrients and leads to high biological productivity \cite{Kennish2016}. As a result, coastal lagoons are considered as natural heritage and classified as areas of wildlife conservation \cite{HERBERT2020631,Nazneen2022} and declared as marine protected areas worldwide due to their rich bioproductivity and valuable environment . 

The transitional nature of coastal lagoons makes them vulnerable to the cumulative effects from the climate fluctuations, specifically rising sea levels \cite{CARRASCO2016356}, flooding \cite{INACIO2023100491}, hydrological disturbances, nutrients availability \cite{LUNARDINI2000135} and human impacts \cite{Luglie}. Richness in natural resources of coastal lagoons attract local population and urges them to actively use resources of the aquatic environment. Hence, coastal lagoon serve as valuable source of food and support the economic development and sustainability for local population which, in turn, leads to the overexploitation of these unique areas and increases anthropogenic pressure \cite{PAULY1994377}. High level of the stress results in the environmental threat to these treasured ecosystems \cite{Davis}. Among recent environmental problems of coastal lagoons are biodiversity and habitat loss \cite{Kumari2022},  organic, chemical and biological pollution including microplastic \cite{GARCESORDONEZ2022120366,BRUSCHI2023110596} and heavy metals \cite{Tripathi2022}, and disrupted land cover patterns due to climate effects. 

\subsection{Theoretical framework}
Monitoring coastal landscapes and variations in land cover types is essential for the management and conservation activities of the Chilika Lake. Such activities are carried out and reported in previous studies based on traditional mapping. Nevertheless, monitoring the land cover types in a lake of such importance using traditional methods is often time-consuming, labour-intensive and includes considerable manual work prone to errors. Though mapping using traditional GIS presents a reliable solution to EO data processing, estimation accuracy is still a notable challenge in mapping coastal areas with high heterogeneity of land cover types. Due to the logical straightforwardness of traditional classification algorithms, their application for thematic GIS-based mapping and analysis of landscape dynamics presents a well-known approach to cartographic workflow with existing case studies on Chilika Lake, Mahanadi Delta and Odisha coastal area \cite{G2019595,MISHRA2023162488,REDDY2014287,HAZRA2022100223}. 

However, spectral complexity of the multispectral satellite images makes recognising land cover types in coastal areas a challenging and less accurate task using k-means clustering or "MaxLike" classification. For instance, specifically for lagoons, optical properties of coastal waters and shelf areas are significantly affected by suspended sediment from the coloured dissolved organic matter (CDOM) which can create noise on EO data. On the other hand, EO data classification using ML methods has been fairly successful. For example, to solve these issues, ML algorithms present automation of image classification \cite{6049966,9297369} which is achieved through computer vision algorithms of pattern recognition and analysis that enables recognising geometrical complexity \cite{990132}. 

Application of ML methods based on Artificial Neural Networks (ANN) applies to Remote Sensing (RS) data processing considerably increases the effectiveness of mapping \cite{rs16010127,rs15041148}. Advanced ML methods enable to automatically reveal spatial and temporal trends landscape dynamics through computer-based algorithms of pattern recognition and data analysis as reported in existing studies \cite{AMANI2023108324,LI2023104653,Lemenkova202021,TSELKA2023131,LATIF202316}. Several ML and AI algorithms exist to analyse and quantify spatial data using analytical and empirical approaches including neural networks which teaches computers to process data in a analytical way which simulates human brain in pattern recognition \cite{HUANG2006489,Rosenblatt}. Among the advanced algorithms of ML are the Random Forest (RF) \cite{598994}, Support Vector Machines (SVM) \cite{Cortes}, Naive Bayes \cite{Zhang2004TheOO}, to mention a few. 

\subsection{Research gap}
There is little research analysing and contrasting landscapes of the Odisha coasts using ML. Existing studies draw generalisations regarding the links between the ecosystems of coastal lagoon of Chilika Lake and the adjacent habitat communities \cite{Sinha2020,BARIK2019352,MISHRA2022150769,BARAL2023103248}. They are however based using the conventional tools of cartographic software which applies traditional mapping methods. To the best of our knowledge, no reported research has been carried out to study the environmental variability of Chilika Lake using ANN techniques. At the same time, scripting libraries are promising tools for cartographic tasks and image processing \cite{9884758,Lemenkova202229,6910592,9553880,Lemenkova202227,ERDEM2021964}. Besides the traditional general-purpose programming languages, diverse modules of GRASS GIS can also be used in image processing \cite{Lemenkova202307}. Hence, ML applications in cartography provide insight to its spatio-temporal variability of landscapes and  environmental processes through classification of the EO data \cite{HAN202387,10292983,990132,WU2024103612}.  

\subsection{Objective and goal}
The main goal of this research is to map and analyse changes of the land cover types surrounding coastal lagoon on the lake using ML learning algorithms of GRASS GIS and EO data. We used Landsat 8-9 OLI/TIRS satellite images on recent six years (2019, 2020, 2021, 2022, 2023 and 2024) processed using ANN method of MLP to map and analyse the spatio-temporal distribution of land cover types around Chilika lagoon, Figure \ref{fig01}. 

To set up the advanced practical background for environmental analysis of Chilika Lake, this study presents the approach of ML for automation of Earth Observation (EO) data processing and visualization. To achieve this goal, the objective is to use the ANN methods by Python's library Scikit-Learn \cite{scikitlearn} which are embedded in the GRASS GIS \cite{GRASS_GIS_software} through ML modules designed for data partition and satellite image processing. The ML approach was selected since it creates a plausible paradigm to map the environmental variability of the coastal landscapes surrounding the lagoon of the Chilika Lake. Among the existing methods, this study employs the MultiLayer Perceptron (MLP) algorithm of ANN which presents an effective solution to image analysis. 

%----------- section ----------->
\section{Study Area}
The coastal lagoon formed by the Chilika Lake lies within the Odisha state, at the estuary of the Daya River, entering the Bay of Bengal on the east coast of India (Figure \ref{fig01}). The area covers spatial extent with the coordinates of 19°28'--19°54' N; 85°06'--85°35' E \cite{PANDA201329}.  The largest brackish water lagoon in Asia \cite{BEHERA2020101328,BARIK201739}, Chilika Lake covers a total area of over $1100 km^2$ with existing fluctuations of the lake surface reported between $1165 km^2$ to $906 km^2$ \cite{Behera2023}. The Chilika Lake is located at the junction of two different water masses -- riverine freshwater and oceanic salt waters from the tidal influx of the Bay of Bengal. These different water fronts interact with the bathymetry of the lagoon and generate local hydrographic setting which vary in salinity, current turbulence and circular flow patterns of waves \cite{Nazneen}. 

A recent geological study has shown that the shallow limnological system of Chilika Lake was a part of the Bay of Bengal during the later stages of the Pleistocene period \cite{BOROLE1993761}, and undergone marked geomorphic evolution during the late Holocene \cite{MISHRA2022457,AMIR2021148}. This included denudation and weathering of the surrounding coasts which is largely caused by climate variability and monsoon effects that influenced the distribution of vegetation types, such as mangroves, within the estuarine environment \cite{DASH2023106754}. Currently, the dynamics of surface area in coastal lagoon of Chilika Lake presents a response to the cumulative effects from tidal morphodynamics, winds and morphometry resulted in fluctuating surface and area which in turn affect the ecosystem of surrounding wetlands.

Coastal lagoon of the Chilika Lake has important conservation features including rare aquatic and sub-aquatic plants, including endemic species, mangrove associations and plants of horticultural importance. Over 300 fish species \cite{MOHAPATRA2015195} and 726 of flowering plants \cite{Sarkar2012} points at its significant biodiversity and ecosystem value. Moreover, Chilika Lake is the largest habitat for migratory waterbirds across India and a home to multiple threatened and rare species of plants and animals \cite{Balachandran2020}. Rich natural resources and unique environment of the Chilika Lake attracted humans to this area since ancient period. Thus, there is archaeological evidence for human settlements existing in the Chilika Lake area which served as marine harbour and port since at least Neolithic period \cite{DASH2022105190}. The occupation of the Chilika Lake is explained by the favourable climate setting and its beneficial location which gave access to the Indian Ocean and ensured maritime trade and international commercial connections in ancient India. The importance of the Chilika Lake in both ecological aspect and for historical value for the development of Indian civilisation resulted in its officially designation a as UNESCO World Heritage site \cite{SETHY2023497} and a Ramsar site \cite{BHUVANAGIRI2018343}.


\begin{figure}[H]
    \begin{center}
%    \setlength{\unitlength}{0.5cm}
	\hspace{-35pt}\includegraphics[width=15.0cm]{Fig_01.jpg}
    \end{center}
    \caption{Topographic map of India with indicated study area showing the location of the Landsat data over the coastal lagoon of Chilika Lake. Software: GMT. Map source: author.}
    \label{fig01}
\end{figure}


%----------- section ----------->
\section{Materials and Methods}

MultiLayer Perceptron (MLP)

%----------- subsection ----------->
\subsection{Data}

\begin{figure}[H]
  \centering
  \begin{tabular}[b]{c}
    \includegraphics[width=6.0cm]{Fig_04a.jpg} \\
    \small (a): 2019 
  \end{tabular}
  \begin{tabular}[b]{c}
    \includegraphics[width=6.0cm]{Fig_04b.jpg} \\
    \small (b): 2020 
  \end{tabular}
    \begin{tabular}[b]{c}
    \includegraphics[width=6.0cm]{Fig_04c.jpg} \\
    \small (c): 2021
  \end{tabular}
  \begin{tabular}[b]{c}
    \includegraphics[width=6.0cm]{Fig_04d.jpg} \\
    \small (d): 2022 
  \end{tabular}
   \begin{tabular}[b]{c}
    \includegraphics[width=6.0cm]{Fig_04e.jpg} \\
    \small (e): 2023 
  \end{tabular}
  \begin{tabular}[b]{c}
    \includegraphics[width=6.0cm]{Fig_04f.jpg} \\
    \small (f): 2024
  \end{tabular}
  \caption{Image processing of the Landsat 8 OLI/TIRS scene on March and November 2014 using unsupervised classification by clustering and Maximal Likelihood method and created colour composites using multispectral bands of Landat images.}
  \label{fig02}
\end{figure}

%----------- subsection ----------->
\subsection{Subsection}

%----------- subsection ----------->
\subsection{Subsection}

%----------- section ----------->
\section{Results}

%----------- section ----------->
\section{Discussion}

%----------- section ----------->
\section{Conclusions}

In this study, we analysed the difference between the effects of diverse ML algorithms on Landsat image analysis and technical relationship of these approaches to image classification. Technically, the study demonstrated the effectiveness of the ML algorithms for image classification which present a powerful alternative to the traditional cartographic tasks through automation. The results of this research show the impact of climate variability on changes in land cover types during short-term time gap and landscape dynamics around the lagoon of the Chilika Lake. The implications of this study are useful for environmental decision makes and supportive for policy makes in sustainable development of Ramsar marine sites of India.

%----------- section ----------->

\institutionalreview{Not applicable.}

\informedconsent{Not applicable.}

\dataavailability{Not applicable.} 

\conflictsofinterest{The author declares no conflict of interest.} 

\abbreviations{Abbreviations}{
The following abbreviations are used in this manuscript:\\

\noindent 
\begin{tabular}{@{}ll}
ANN & Artificial Neural Networks \\
CDOM & Coloured Dissolved Organic Matter \\
DL & Deep Learning \\
EO & Earth Observation \\
GEBCO & General Bathymetric Chart of the Oceans \\
GMT & Generic Mapping Tools \\
GRASS & Geographic Resources Analysis Support System \\
GIS & Geographic Information System \\
Landsat OLI/TIRS & Landsat Operational Land Imager and Thermal Infrared Sensor \\
ML & Machine Learning \\
MLP & MultiLayer Perceptron \\
RF & Random Forest \\
RS & Remote Sensing \\
SVM & Support Vector machines \\
USGS & United States Geological Survey \\
UTM & Universal Transverse Mercator \\
\end{tabular}
}

%%%%%%%%%%%%%%%%%%%%%%%%%%%%%%%%%%%%%%%%%%
%% Optional
\appendixtitles{no} % Leave argument "no" if all appendix headings stay EMPTY (then no dot is printed after "Appendix A"). If the appendix sections contain a heading then change the argument to "yes".
\appendixstart
\appendix
\section[\appendixname~\thesection]{}
\subsection[\appendixname~\thesubsection]{}
The appendix 

\section[\appendixname~\thesection]{}
All appendix sections must be cited in the main text. In the appendices, Figures, Tables, etc. should be labeled, starting with ``A''---e.g., Figure A1, Figure A2, etc.

%%%%%%%%%%%%%%%%%%%%%%%%%%%%%%%%%%%%%%%%%%
\begin{adjustwidth}{-\extralength}{0cm}
%\printendnotes[custom] % Un-comment to print a list of endnotes

\reftitle{References}
\nocite{*}

%=====================================
% References, variant A: external bibliography
%=====================================
\bibliography{Chilika.bib}

%%%%%%%%%%%%%%%%%%%%%%%%%%%%%%%%%%%%%%%%%%
\end{adjustwidth}
\end{document}
